\begin{center}
\textbf{{\huge 5.Conducting the experiment}}
\end{center}
\vspace{ 2 mm}
At first,  We took just one rubber band and decided to use that as the representative rubber band of the population of rubber bands. We took several measurements using the  rubber band. After gathering several data points we saw that the rubber band was no longer useful to gather data. This led us to conclude that only one rubber band can not used to gather data. Because as time passes and we conduct experiments using a rubber band , the rubber band suffers wear and tear. This leads to error that can be controlled if we use several rubber bands to  gather data. So, this time we decided to use 9 rubber bands selected randomly  from the bunch of rubber bands. We decided to use each rubber band for a set number of times to gather data. At specific pulled length we took two measurements for each rubber band. Then we averaged the data to associate that average as travelled distance with that specific stretched length. We did this close to 180 times. We  took data in a manner so that we can identify which data came from which rubber band. After completion of data gathering stage we did some computations on the data to calculate actual stretched length. We plotted the data using R. For most of the rubber bands the corresponding plot suggested linear relation. This suggested us to use linear regression to model it. As we wanted to use our model for any rubber band from the population we used, We mixed the data and plotted it. The data cloud still contained the essential feature of being linear like. So seeing this we went ahead and used linear regression.
